% Software and the structure of the notebook of the Day 6 lab.
\begin{frame}{Software and files}
  \begin{columns}[T]
    \begin{column}{0.47\textwidth}
      \textbf{You need}
      \begin{itemize}
        \item A laptop. No board.
        \item The Python environment of Days 1 and 3, with
              \code{tensorflow-cpu} and \code{ai-edge-litert}.
        \item No new package.
      \end{itemize}
    \end{column}
    \begin{column}{0.49\textwidth}
      \textbf{The lab folder gives}
      \begin{itemize}
        \item \code{compression.ipynb}: the notebook with seven tasks.
        \item \code{models/baseline.keras}: the CNN of the lecture.
        \item \code{models/teacher.keras}: a wider CNN.
      \end{itemize}
    \end{column}
  \end{columns}
  \vskip 0.6em
  \small
  \renewcommand{\arraystretch}{1.15}
  \begin{tabular}{@{}lL{0.80\textwidth}@{}}
    Lab files & \code{Labs/day06/}. Run the commands from this folder. \\
    Notebook  & \code{../../.venv/bin/jupyter lab compression.ipynb} \\
    Dataset   & Fashion-MNIST. The first cell downloads it: about 30 MB. \\
  \end{tabular}
  \vskip 0.4em
  \begin{itemize}
    \item The two models are trained. A training from the start needs too
          much time for a lab.
  \end{itemize}
\end{frame}

\begin{frame}{From two models to one plot}
  \begingroup\centering
  \begin{tikzpicture}[font=\footnotesize, >={Stealth[length=5pt]},
      b/.style={draw=coolgray, rounded corners=3pt, fill=boxgray,
                minimum width=2.2cm, minimum height=0.9cm, align=center},
      lab/.style={font=\scriptsize, text=coolgray, align=center}]
    \node[b] (base) at (0,0.7)  {Baseline};
    \node[b] (teach) at (0,-0.9) {Teacher};
    \node[b, draw=darkcyan] (prune) at (2.9,0.7) {Part A:\\pruning};
    \node[b, draw=treegreen] (dist) at (2.9,-0.9) {Part B:\\distillation};
    \node[b, draw=cardinalred] (quant) at (5.8,-0.1) {Part C:\\quantize};
    \node[b] (plot) at (8.5,-0.1) {Plot and\\selection};
    \draw[->] (base) -- (prune); \draw[->] (teach) -- (dist);
    \draw[->] (prune.east) -- (quant.west);
    \draw[->] (dist.east) -- (quant.west);
    \draw[->] (quant) -- (plot);
    \node[lab] at (2.9,1.45) {4 candidates};
    \node[lab] at (2.9,-1.65) {2 candidates};
    \node[lab] at (5.8,-0.85) {7 models,\\14 files};
  \end{tikzpicture}\par\endgroup
  \vskip 0.4em
  \small
  \begin{itemize}
    \item Each part adds models to one list of candidates. The baseline is
          the first candidate.
    \item Part C converts each candidate to a \code{float32} file and to an
          \code{int8} file, and it measures four numbers for each file: the
          size, the accuracy, the latency, and the peak of the activations.
    \item An experiment starts only when the check of its task prints
          \code{complete}.
  \end{itemize}
  \keypoint{The latency of this lab is the latency of your laptop. Day 9
    measures the same models on the boards.}
\end{frame}
